\documentclass[10pt,a4paper]{amsart}
\usepackage[margin=19mm]{geometry}
\usepackage{amsmath,amssymb,mathtools}
\usepackage[scr=rsfs,cal=euler]{mathalfa}
\usepackage{xfrac}
\usepackage[unicode,hidelinks,bookmarks=false]{hyperref}
\newcommand{\ie}{i.e.\ }
\newcommand{\cf}{cf.\ }
\newcommand{\resp}{respectively\ }
\newcommand{\Subsection}{\subsection}
\providecommand{\nobreakdash}{\nobreak\hbox{-}}
\setlength{\emergencystretch}{2em}
\multlinegap0pt
\usepackage{mathtools}
\usepackage{amssymb}
\usepackage[shortlabels]{enumitem}
\newtheorem{lemma}{Lemma}
\title{Spinor inequality for magnetic fields on spin manifolds}
\author{Jurgen Julio-Batalla}
\date{Source: arXiv:2603.23218v1 (CC BY 4.0)}
\begin{document}
\maketitle

\noindent\emph{Source and licence.} Spinor inequality for magnetic fields on spin manifolds. Version arXiv:2603.23218v1 (CC BY 4.0), distributed by arXiv under the Creative Commons Attribution 4.0 International licence, http://creativecommons.org/licenses/by/4.0/, as recorded in the arXiv licence metadata for this exact version and shown on the abstract page https://arxiv.org/abs/2603.23218v1. Full licence notice: \texttt{licenses/arxiv-2603.23218-CC-BY-4.0.txt}.\par\smallskip

\noindent\textit{Editorial scope.} Counted unit: the article's lemma lemma (source0.tex lines 274--276). The article's complete original proof of it is delivered with it (source0.tex lines 277--304, one \texttt{proof} environment of 28 lines). No mathematics is written, repaired or re-derived: the statement and the proof are the article's own lines, in the article's own notation, and no retained line is substituted or re-flowed.  No further source-local context is needed: every cross-reference of every retained line resolves inside the two delivered spans.  Nothing was omitted from any retained span, so the package carries no omission record.  No retained line contains a citation, so the wrapper carries no bibliography and no bibliography entry.  No retained line contains a figure environment, an image-inclusion command or drawing code, so no image is delivered and the package declares no asset dependency.  Editorial formatting only, in addition: an AMS \texttt{amsart} preamble that restates the article's own declarations and macros used by the retained lines, namely the environments \texttt{lemma} on separate counters, together with the standard packages those lines need.  The retained environments print in this document as Lemma 1 in order of appearance, whereas the article prints the same items with the numbering of its own full document.  The complete native source of the article as distributed in the arXiv e-print for this exact version is bundled unchanged as \texttt{sources/197063/source0.tex}.
\begin{lemma}
The quantity $$\lambda_1\left(4i\frac{\langle dA\cdot_{g_0}\varphi,\varphi\rangle_{g_0}}{|\varphi|_{g_0}^2},g_0\right)$$ is conformally invariant.
\end{lemma}

\begin{proof}
Let $g=h^{4/(n-2)}g_0$ and let $F$ be the canonical isomorphism from $S(M,g_0,\sigma)$ to $S(m,g,\sigma)$ from Section 2.
We will prove that $\lambda_1$ with respect to the metric $g_0$ equals that with respect to the metric $g$.

Assume $(\varphi,A)$ is a nontrivial solution of the zero mode equation for the metric $g_0$. As in Section 2, the pair $(\psi,A)$ defined as $\psi=h^{-(n-1)/(n-2)}F(\varphi)$ is a solution of the equation with respect to the metric $g$.

We define the functional
$$I_{(g,\psi,A)}(u)=\dfrac{\int_M(a_n|\nabla u|_g^2+s_gu^2)dv_g}{4i\int_M \frac{\langle dA\cdot_g\psi,\psi\rangle_g}{|\psi|_g^2}  u^2dv_g  }.$$

First note that by the conformal transformation of the conformal Laplacian $$\int_M(a_n|\nabla u|_g^2+s_gu^2)dv_g=\int_M L_g(u)u\;dv_g=\int_Mh^{-(n+2)/(n-2)}L_{g_0}(hu)u\;dv_{g}.$$
Hence $$\int_M(a_n|\nabla u|_g^2+s_gu^2)dv_g=\int_ML_{g_0}(hu)hu\;dv_{g_0}$$
Second, $$|\psi|_{g}^2=h^{-2(n-1)/(n-2)}|\varphi|_{g_0}^2.$$

For a local $g_0-$orthonormal frame $\{E_j\}$  we consider the corresponding local $g-$orthonormal frame $\{\bar{E_j}:=h^{-2/(n-2)}E_j\}.$

Thus 
$$
dA\cdot_{g}\psi=\sum\limits_{j<k}dA(\bar{E_j},\bar{E_k})\bar{E_j}\cdot_{g} \bar{E_k}\cdot_{g}\psi.
$$
Since $$\langle \bar{E_j}\cdot_{g} \bar{E_k}\cdot_{g}F(\varphi),F(\varphi)\rangle_g=\langle E_j\cdot_{g_0} E_k\cdot_{g_0}\varphi,\varphi\rangle_{g_0}$$
we obtain $$\langle dA\cdot_{g}\psi,\psi\rangle_g=h^{-2(n-1)/(n-2)}\sum\limits_{j<k}dA(\bar{E_j},\bar{E_k})\langle E_j\cdot_{g_0} E_k\cdot_{g_0}\varphi,\varphi\rangle_{g_0}.$$
Then $$\langle dA\cdot_{g}\psi,\psi\rangle_g=h^{\frac{-4-2(n-1)}{n-2}}\langle dA\cdot_{g_0}\varphi,\varphi\rangle_{g_0}.$$
From this, we obtain 
$$\int_M \frac{\langle dA\cdot_g\psi,\psi\rangle_g}{|\psi|_g^2}  u^2dv_g =\int_M h^{-4/(n-2)}\frac{\langle dA\cdot_{g_0}\varphi,\varphi\rangle_{g_0}}{|\varphi|_{g_0}^2}  udv_{g}=\int_M \frac{\langle dA\cdot_{g_0}\varphi,\varphi\rangle_{g_0}}{|\varphi|_{g_0}^2}  (hu)^2dv_{g_0} $$
Therefore $$I_{(g,\psi,A)}(u)=I_{(g_0,\varphi,A)}(hu).$$

The lemma follows by taking the infimum in the previous equality. 
\end{proof}

\end{document}
